\documentclass[11pt]{article}

\usepackage[margin=1in]{geometry}
\usepackage{amsmath}
\usepackage{amssymb}
\usepackage{hyperref}

\title{Mathematical Foundations \\ \large Emission-Trajectory ML Project}
\author{Dr. Khawar Naeem \\ Qatar Transportation and Traffic Safety Center (QTTSC), Qatar University}
\date{14 July 2026 (living document; grows with each modeling session)}

\begin{document}
\maketitle

\begin{abstract}
This document formalizes the mathematics behind the emission-trajectory
forecasting pipeline (\url{https://github.com/khawar89/energy-carbon-forecast-ml}):
the evaluation metrics every model is scored by, the two baselines that
define the bar a model must clear, Ridge regression as the project's first
learned model, the ensemble and boosting methods that follow it, and the
uncertainty quantification that decides whether a skill difference is
statistically real. Each
equation is numbered and tied directly to the corresponding code in
\texttt{src/evaluate.py}, \texttt{notebooks/03\_baselines.ipynb},
\texttt{notebooks/04\_models\_ridge\_tree.ipynb}, and
\texttt{notebooks/05\_xgboost.ipynb}. New sections are added as
later sessions (error analysis) are built, so this remains a
growing companion to the coding sessions rather than a one-time summary.
This is intended to support future methodology writing and research
extension of the project, not to substitute for it.
\end{abstract}

\tableofcontents

\section{Notation}
\label{sec:notation}

This document uses consistent notation across all sections, matching the
pipeline built in \texttt{src/build\_features.py} and \texttt{src/evaluate.py}.

\begin{itemize}
  \item $c$ indexes a country; $t$ indexes a calendar year.
  \item $y_{c,t}$ denotes the observed production-based CO$_2$ emissions
    (Mt) of country $c$ in year $t$.
  \item The \emph{feature year} $t$ is the year whose information is known
    at prediction time; the \emph{target year} $t+1$ is the year being
    predicted. No feature used to predict $y_{c,t+1}$ may use information
    from after year $t$ (see the project's \texttt{CLAUDE.md} and
    \texttt{AGENTS.md} for the full leakage-prevention rationale; every
    equation in this document assumes that constraint is already
    enforced).
  \item $\hat y_{c,t+1}$ denotes a model's prediction of $y_{c,t+1}$, made
    using information available through year $t$.
  \item $x_{c,t} \in \mathbb{R}^p$ denotes the feature vector for country
    $c$ at year $t$ (lags, rolling statistics, population, and related
    columns), with $p$ the number of features.
  \item $n$ denotes the number of (country, year) rows in a given
    evaluation set (for example, the validation split).
  \item Splits are labeled train, val (validation), and test, assigned by
    the \emph{target} year $t+1$, not the feature year $t$: train covers
    target years through 2014, val covers 2015--2018, test covers
    2019--2023, and a separate provisional table covers target year 2024.
\end{itemize}

\section{Evaluation Metrics}
\label{sec:metrics}

Every model in this project, from the persistence baseline through
XGBoost, is scored by the same set of metrics, implemented once in
\texttt{src/evaluate.py} and reused unchanged in every later session. This
section states each metric formally.

\subsection{Mean Absolute Error (MAE)}

\begin{equation}
\label{eq:mae}
\mathrm{MAE} = \frac{1}{n}\sum_{i=1}^{n} \left| y_i - \hat y_i \right|
\end{equation}

MAE is the average size of a miss, in the original units (Mt CO$_2$). It is
interpretable, but because errors are measured in absolute units, its
value is dominated by the largest countries by emissions level.

\subsection{Root Mean Squared Error (RMSE)}

\begin{equation}
\label{eq:rmse}
\mathrm{RMSE} = \sqrt{\frac{1}{n}\sum_{i=1}^n \left(y_i - \hat y_i\right)^2}
\end{equation}

Squaring each error before averaging penalizes large misses quadratically.
Consequently $\mathrm{RMSE} \gg \mathrm{MAE}$ for a given model signals
that its error budget is concentrated in a small number of large misses (a
``fat tail'') rather than spread evenly across rows.

\subsection{Median Absolute Error (MedianAE)}

\begin{equation}
\label{eq:medianae}
\mathrm{MedianAE} = \operatorname{median}_i \left| y_i - \hat y_i \right|
\end{equation}

MedianAE is immune to the largest emitters: it reports the typical error
for the typical country. Comparing Equations~\ref{eq:mae}
and~\ref{eq:medianae} on the same model reveals how much of the average
error is attributable to a handful of very large countries rather than to
the rest of the sample, a scale-dependence issue discussed generally by
\cite{hyndman2006another}.

\subsection{Median Absolute Percentage Error (MdAPE)}

\begin{equation}
\label{eq:mdape}
\mathrm{MdAPE} = 100 \times \operatorname{median}_{i \,:\, y_i \ge \tau}
\left( \frac{|y_i - \hat y_i|}{y_i} \right)
\end{equation}

where $\tau$ is a minimum-level threshold (this project uses
$\tau = 1$~Mt). Percentage errors are undefined or explosive as
$y_i \to 0$; the exclusion in Equation~\ref{eq:mdape} is part of the
metric's definition, not an afterthought, following the general caution
around percentage-based accuracy measures raised by
\cite{hyndman2006another}.

\subsection{Persistence Skill Score}

\begin{equation}
\label{eq:skill}
\mathrm{Skill} = 1 - \frac{\mathrm{MAE}_{\mathrm{model}}}{\mathrm{MAE}_{\mathrm{persistence}}}
\end{equation}

The skill score in Equation~\ref{eq:skill} measures a model's error
relative to the persistence baseline (Section~\ref{sec:baselines}), rather
than in absolute terms. $\mathrm{Skill}=0$ means the model is exactly as
accurate as assuming no change; $\mathrm{Skill}=1$ is the unreachable
limit of zero error; negative skill means the model is less accurate than
doing nothing and must be reported as a loss, not merely a smaller win.

\subsection{The same-rows rule}

Every comparison table in this project scores all candidate models on an
identical set of rows: a row missing any model's prediction, or missing
the target itself, is dropped for every model before
Equations~\ref{eq:mae}--\ref{eq:skill} are computed
(\texttt{evaluate.comparison\_table}). Without this rule, a model could
appear to win only because it was scored on an easier subset of rows.

\section{Baselines}
\label{sec:baselines}

A baseline is a model that requires no fitting: no parameters are
estimated from data. Its purpose is to define the bar every learned model
must clear (Section~\ref{sec:metrics}'s skill score is measured against
it).

\subsection{Persistence}

\begin{equation}
\label{eq:persistence}
\hat y_{c,t+1} = y_{c,t}
\end{equation}

Equation~\ref{eq:persistence} predicts that next year's level equals this
year's level. This is the naive forecast, and it is the optimal forecast
under a random-walk-without-drift model of the underlying series
\cite{hyndman2021forecasting}. It is a strong baseline for country-level
CO$_2$ specifically because the underlying physical stock (power plants,
vehicle fleets, industrial capital) changes slowly: the median absolute
year-over-year change in this dataset since 2010 is approximately 4.4\%
of the level (Session 1 EDA).

\subsection{Linear Trend (Drift)}

\begin{equation}
\label{eq:trend}
\hat y_{c,t+1} = y_{c,t} + \bar\Delta_{c,t}^{(5)},
\qquad
\bar\Delta_{c,t}^{(5)} = \frac{y_{c,t} - y_{c,t-5}}{5}
\end{equation}

Equation~\ref{eq:trend} extrapolates the average yearly change over the
preceding five years one step forward. It is this project's
five-year-windowed variant of the classical drift method
\cite{hyndman2021forecasting}, which in its general form draws a line
between a series' first and last observed points and extrapolates it. The
trade-off is explicit: Equation~\ref{eq:trend} outperforms
Equation~\ref{eq:persistence} for a country on a sustained trajectory, and
underperforms it the moment a country's direction reverses, since the
linear extrapolation has no mechanism to anticipate a turn.

\subsection{Bias direction}

For any series with a genuine sustained trend, persistence
(Equation~\ref{eq:persistence}) is a biased estimator of $y_{c,t+1}$: if
$y_{c,t}$ is rising, persistence systematically under-predicts; if
$y_{c,t}$ is falling (for example, a country whose emissions peaked and
have since declined), persistence systematically over-predicts. The sign
of the bias tracks the sign of the country's own recent trend, not a
fixed direction.

\section{Ridge Regression}
\label{sec:ridge}

\subsection{Ordinary least squares}

Given a design matrix $X \in \mathbb{R}^{n\times p}$ (rows are
(country, year) observations, columns are features) and target vector
$y \in \mathbb{R}^n$, ordinary least squares (OLS) solves

\begin{equation}
\label{eq:ols}
\hat\beta_{\mathrm{OLS}} = \arg\min_{\beta} \left\| y - X\beta \right\|_2^2
\end{equation}

with closed-form solution
$\hat\beta_{\mathrm{OLS}} = (X^\top X)^{-1}X^\top y$, provided
$X^\top X$ is invertible.

\subsection{The problem with correlated features}

This project's feature set includes multiple lags and rolling statistics
of the same underlying series (\texttt{co2\_lag1}, \texttt{co2\_lag3},
\texttt{co2\_lag5}, \texttt{co2\_lag10}, \texttt{co2\_roll5\_mean},
\texttt{co2\_roll10\_mean}), which are, by construction, highly correlated
with one another. When columns of $X$ are strongly correlated,
$X^\top X$ is close to singular, and $\hat\beta_{\mathrm{OLS}}$ becomes
highly sensitive to small changes in the data: its variance inflates even
though it remains unbiased.

\subsection{Ridge regression}

Ridge regression adds an $\ell_2$ penalty on the coefficients:

\begin{equation}
\label{eq:ridge}
\hat\beta_{\mathrm{ridge}} = \arg\min_{\beta}
\left\| y - X\beta \right\|_2^2 + \alpha \left\| \beta \right\|_2^2,
\qquad \alpha \ge 0
\end{equation}

with closed-form solution

\begin{equation}
\label{eq:ridge-closed-form}
\hat\beta_{\mathrm{ridge}} = \left(X^\top X + \alpha I\right)^{-1} X^\top y
\end{equation}

Adding $\alpha I$ to $X^\top X$ in Equation~\ref{eq:ridge-closed-form}
guarantees invertibility for any $\alpha>0$ and shrinks the coefficients of
correlated features toward one another and toward zero. This trades a
small amount of bias for a reduction in variance \cite{hoerl1970ridge},
and is the reason this project prefers Ridge over OLS given its
correlated lag and rolling-statistic feature set.

\subsection{Preprocessing}

Because the ridge penalty in Equation~\ref{eq:ridge} is applied uniformly
to all coefficients, features must be on comparable scales before
fitting; unlike tree-based models (Section~\ref{sec:ensembles}), Ridge
requires feature scaling, and that scaler must be fit on training rows
only, to avoid leaking validation or test distributional information into
its parameters.

\section{Ensembles and Boosting}
\label{sec:ensembles}

\subsection{Regression trees}

A regression tree partitions the feature space into disjoint regions
$R_1,\dots,R_J$ and predicts the mean target value within each region. A
single split at feature $j$, threshold $s$ is chosen to minimize the total
sum of squared residuals across the two resulting regions:

\begin{equation}
\label{eq:tree-split}
(j^*, s^*) = \arg\min_{j,s}
\left[
  \sum_{i:\, x_i \in R_1(j,s)} \left(y_i - \bar y_{R_1}\right)^2
  +
  \sum_{i:\, x_i \in R_2(j,s)} \left(y_i - \bar y_{R_2}\right)^2
\right]
\end{equation}

Trees require no feature scaling: Equation~\ref{eq:tree-split} depends
only on the ordering of feature values within a column, not their
magnitude, unlike the Ridge objective in Equation~\ref{eq:ridge}.

\subsection{Bagging}

Bagging (bootstrap aggregating) reduces variance by averaging many trees,
each fit to an independent bootstrap resample of the training rows:

\begin{equation}
\label{eq:bagging}
\hat f_{\mathrm{bag}}(x) = \frac{1}{B}\sum_{b=1}^{B} \hat f_b(x)
\end{equation}

where each $\hat f_b$ is a tree fit to bootstrap sample $b$
\cite{breiman1996bagging}. Random forests extend
Equation~\ref{eq:bagging} by additionally restricting each split in
Equation~\ref{eq:tree-split} to a random subset of features, further
decorrelating the individual trees \cite{breiman2001random}.

\subsection{Gradient boosting}

Where bagging averages independent trees fit in parallel, gradient
boosting fits trees sequentially, each one correcting the errors of the
ensemble built so far. The model is an additive expansion

\begin{equation}
\label{eq:gb-additive}
F_M(x) = \sum_{m=1}^{M} \gamma_m h_m(x)
\end{equation}

built greedily: at each stage $m$, tree $h_m$ is fit to the negative
gradient (pseudo-residual) of the loss function $L$ with respect to the
current model's predictions,

\begin{equation}
\label{eq:gb-pseudoresidual}
r_{i,m} = -\left.
\frac{\partial L(y_i, F(x_i))}{\partial F(x_i)}
\right|_{F=F_{m-1}}
\end{equation}

which for squared-error loss reduces to the ordinary residual
$y_i - F_{m-1}(x_i)$. This is gradient descent carried out in the space of
functions rather than parameters \cite{friedman2001greedy}.
\texttt{HistGradientBoostingRegressor} (scikit-learn), used in this
project's Session 4, implements
Equations~\ref{eq:gb-additive}--\ref{eq:gb-pseudoresidual} with
histogram-binned feature values, which reduces the split-search cost of
Equation~\ref{eq:tree-split} from scanning every unique value to scanning
a fixed number of bins, and handles missing feature values natively.

\subsection{XGBoost}
\label{sec:xgboost}

XGBoost augments gradient boosting with an explicit regularization term
and a second-order (Newton) approximation of the loss. Its objective at
boosting round $m$ is

\begin{equation}
\label{eq:xgb-objective}
\mathcal{L}^{(m)} =
\sum_{i=1}^n l\!\left(y_i,\, F_{m-1}(x_i) + h_m(x_i)\right)
+ \Omega(h_m),
\qquad
\Omega(h) = \gamma T + \tfrac{1}{2}\lambda \|w\|^2
\end{equation}

where $T$ is the number of leaves in tree $h_m$ and $w$ its leaf weights.
This project treats XGBoost as a comparison model, not the project's goal
(see \texttt{AGENTS.md}), with a small, defensible hyperparameter grid
selected on the validation years rather than a wide search, implemented
in \texttt{notebooks/05\_xgboost.ipynb} and \texttt{src/train.py}.

\subsubsection{The second-order (Newton) step}

Rather than fitting $h_m$ to the first-order pseudo-residual alone
(Equation~\ref{eq:gb-pseudoresidual}), XGBoost approximates the loss in
Equation~\ref{eq:xgb-objective} by its second-order Taylor expansion
around the current prediction $F_{m-1}(x_i)$. Writing, for each
observation,

\begin{equation}
\label{eq:xgb-grad-hess}
g_i = \left.\frac{\partial\, l(y_i, F)}{\partial F}\right|_{F=F_{m-1}(x_i)},
\qquad
h_i = \left.\frac{\partial^2 l(y_i, F)}{\partial F^2}\right|_{F=F_{m-1}(x_i)},
\end{equation}

(the Hessian $h_i$ carries an observation subscript $i$; the tree $h_m$
carries a round subscript $m$, so the two symbols do not collide), the
objective becomes, up to a constant that does not depend on $h_m$,

\begin{equation}
\label{eq:xgb-taylor}
\tilde{\mathcal{L}}^{(m)} =
\sum_{i=1}^{n}\left[ g_i\, h_m(x_i) + \tfrac{1}{2} h_i\, h_m(x_i)^2 \right]
+ \gamma T + \tfrac{1}{2}\lambda \sum_{j=1}^{T} w_j^2 .
\end{equation}

A tree with fixed structure assigns every observation to exactly one leaf
$j$ with constant value $w_j$. Collecting the per-leaf sums
$G_j = \sum_{i \in I_j} g_i$ and $H_j = \sum_{i \in I_j} h_i$, where
$I_j$ is the set of observations in leaf $j$, Equation~\ref{eq:xgb-taylor}
separates across leaves:

\begin{equation}
\label{eq:xgb-per-leaf}
\tilde{\mathcal{L}}^{(m)} =
\sum_{j=1}^{T}\left[ G_j w_j + \tfrac{1}{2}\left(H_j + \lambda\right) w_j^2 \right]
+ \gamma T .
\end{equation}

Each bracket is a one-dimensional quadratic in $w_j$, minimized in closed
form:

\begin{equation}
\label{eq:xgb-leaf-weight}
w_j^{*} = -\frac{G_j}{H_j + \lambda},
\qquad
\tilde{\mathcal{L}}^{(m)}_{\min} =
-\tfrac{1}{2}\sum_{j=1}^{T}\frac{G_j^2}{H_j + \lambda} + \gamma T .
\end{equation}

\subsubsection{Split gain and pruning}

Because Equation~\ref{eq:xgb-leaf-weight} scores any fixed tree structure,
the value of splitting one leaf into left and right children ($L$, $R$) is
the reduction in the minimized objective, the split gain
\cite{chen2016xgboost}:

\begin{equation}
\label{eq:xgb-gain}
\mathrm{Gain} = \tfrac{1}{2}\left[
\frac{G_L^2}{H_L + \lambda} + \frac{G_R^2}{H_R + \lambda}
- \frac{(G_L + G_R)^2}{H_L + H_R + \lambda}
\right] - \gamma .
\end{equation}

Splits are grown greedily by maximizing Equation~\ref{eq:xgb-gain} over
candidate features and thresholds; any split whose gain is negative after
the $\gamma$ deduction is not made, so $\gamma$ acts as a pruning
threshold on weak splits. (The xgboost implementation reports gain without
the constant factor $\tfrac{1}{2}$, which does not change which split is
chosen; verified against xgboost 3.3.0, whose reported gain is exactly
twice Equation~\ref{eq:xgb-gain} at $\gamma = 0$.)

\subsubsection{Specialization to squared error, and what $\lambda$ does}

Under the squared-error loss $l = \tfrac{1}{2}(y_i - F)^2$ used in this
project, Equation~\ref{eq:xgb-grad-hess} gives
$g_i = F_{m-1}(x_i) - y_i = -r_i$ (the negative residual) and $h_i = 1$,
so for a leaf containing $n_j$ observations Equation~\ref{eq:xgb-leaf-weight}
becomes

\begin{equation}
\label{eq:xgb-shrunk-leaf}
w_j^{*} = \frac{\sum_{i \in I_j} r_i}{n_j + \lambda}
        = \bar r_j \cdot \frac{n_j}{n_j + \lambda},
\end{equation}

where $\bar r_j$ is the mean residual in the leaf. Plain gradient boosting
(Section on gradient boosting above) would output $\bar r_j$ itself;
XGBoost multiplies it by $n_j / (n_j + \lambda)$, shrinking every leaf's
correction toward zero and shrinking thin-evidence leaves (small $n_j$)
hardest. This is the second-order machinery's practical effect on this
project's data: corrections supported by few country-year rows are damped
before they can overfit. Equations~\ref{eq:xgb-leaf-weight},
\ref{eq:xgb-gain}, and~\ref{eq:xgb-shrunk-leaf} were verified numerically
against xgboost 3.3.0 on a worked four-observation example (a single
depth-one tree, $\lambda = 1$): the fitted leaf values equal both
closed-form expressions exactly, and the split chosen is the gain-maximal
one.

\subsubsection{Shrinkage, subsampling, and early stopping}

Each fitted tree enters the ensemble scaled by a learning rate
$\eta \in (0, 1]$:

\begin{equation}
\label{eq:xgb-shrinkage}
F_m(x) = F_{m-1}(x) + \eta\, h_m(x),
\end{equation}

so a smaller $\eta$ requires more boosting rounds $M$ to reach the same
fit; $\eta$ and $M$ trade off directly, and slow learning with more trees
generalizes better at higher compute cost \cite{friedman2001greedy}.
XGBoost additionally fits each tree on a random fraction of rows
(\texttt{subsample}) and columns (\texttt{colsample\_bytree}), which
decorrelates successive trees in the same spirit as bagging
\cite{chen2016xgboost}.

The number of rounds $M$ is not tuned by grid search in this project but
chosen by early stopping: trees are added while the validation-set MAE
improves, and boosting stops when it has not improved for a fixed patience
window, freezing

\begin{equation}
\label{eq:xgb-early-stop}
M^{*} = \arg\min_{m \le M_{\max}} \mathrm{MAE}_{\mathrm{val}}(F_m).
\end{equation}

Early stopping is model selection on the validation years, which this
project's protocol allows; selecting anything by test-set performance is
not. For the single test evaluation, each model is refit on the pooled
training and validation rows (target years through 2018) with the tree
count fixed at $M^{*}$ and early stopping disabled, so the test years
never influence any modeling choice, including the number of trees.

\subsubsection{Why trees require the delta parameterization}

A regression tree predicts a constant within each region of
Equation~\ref{eq:tree-split}, so an ensemble of trees can never predict a
level outside the range of its training targets. A country whose emissions
rise beyond anything seen in training saturates at the training ceiling
and is systematically under-predicted in the level parameterization of
Equation~\ref{eq:level-delta}. Predicting the bounded year-over-year
change $\delta_{c,t+1}$ and reconstructing the level from the known
current value sidesteps this failure mode, which is why the tree-based
models in this project are evaluated in both parameterizations.

\section{Uncertainty and Significance of a Skill Score}
\label{sec:uncertainty}

The metrics of Section~\ref{sec:metrics} are point estimates. A skill
score of, say, $-0.078$ reports that a model's average error is 7.8 percent
larger than persistence's on the evaluation sample, but it does not say
whether that gap is distinguishable from zero or is within sampling noise.
For a comparison to support a claim, the gap needs an interval and a
significance statement. This section formalizes both, implemented as
optional functions in \texttt{src/evaluate.py} that are deliberately not
wired into the mandatory comparison table.

\subsection{Loss differential}

Let $e_i^{\mathrm{m}} = |y_i - \hat y_i^{\mathrm{m}}|$ and
$e_i^{\mathrm{p}} = |y_i - \hat y_i^{\mathrm{p}}|$ be the absolute errors of
a candidate model and of persistence on row $i$. Their difference

\begin{equation}
\label{eq:loss-diff}
d_i = e_i^{\mathrm{p}} - e_i^{\mathrm{m}}
\end{equation}

is positive when the model beats persistence on row $i$. The classical
Diebold--Mariano test \cite{diebold1995comparing} assesses the null
hypothesis $\mathbb{E}[d_i]=0$ (equal predictive accuracy) using the mean
and a serial-correlation-robust variance of $\{d_i\}$. The skill score of
Equation~\ref{eq:skill} is a monotone reparameterization of the ratio of
mean absolute errors, so testing whether skill differs from zero is the
same question as testing whether the mean loss differential
(Equation~\ref{eq:loss-diff}) differs from zero.

\subsection{Why a row-level interval is wrong here}

The Diebold--Mariano test was designed for a single time series of
forecasts. This project's evaluation sample is a panel: each country
contributes several years, and a country's own years are correlated
(its emissions, and its errors, move together). Treating the $n$ rows as
independent would understate the sampling variance, because the effective
number of independent observations is closer to the number of countries
than to the number of rows. In addition, a few large emitters dominate the
absolute error (Section~\ref{sec:metrics}), so the influence of individual
countries is uneven. Both features call for resampling at the level of the
country, not the row.

\subsection{Country-clustered bootstrap interval}

Let $\mathcal{C}$ be the set of countries in the evaluation sample, with
$|\mathcal{C}| = C$. A cluster bootstrap \cite{cameron2008bootstrap},
\cite{efron1993bootstrap} draws $C$ countries with replacement, pools all
rows of the drawn countries, and recomputes the statistic. For bootstrap
replicate $b = 1,\dots,B$, let $\mathcal{C}^{(b)}$ be the multiset of drawn
countries and $R^{(b)}$ the pooled rows; the resampled skill is

\begin{equation}
\label{eq:boot-skill}
s^{(b)} = 1 - \frac{\sum_{i \in R^{(b)}} e_i^{\mathrm{m}}}
                   {\sum_{i \in R^{(b)}} e_i^{\mathrm{p}}}
\end{equation}

Resampling whole countries preserves the within-country dependence that a
row-level resample would destroy. A $100(1-\alpha)$ percent confidence
interval for the skill is the empirical interval of the replicates,

\begin{equation}
\label{eq:boot-ci}
\left[\; q_{\alpha/2}\!\left(\{s^{(b)}\}\right),\;
        q_{1-\alpha/2}\!\left(\{s^{(b)}\}\right) \;\right]
\end{equation}

where $q_\tau$ is the $\tau$ empirical quantile. A two-sided bootstrap
$p$-value for the null hypothesis that the model ties persistence
($s = 0$) is

\begin{equation}
\label{eq:boot-p}
p = 2 \min\!\left(
  \frac{1}{B}\sum_{b=1}^{B}\mathbf{1}\{s^{(b)} \le 0\},\;
  \frac{1}{B}\sum_{b=1}^{B}\mathbf{1}\{s^{(b)} \ge 0\}
\right),
\end{equation}

clipped to $[0,1]$. This is a bootstrap tail probability, not an exact
analytic test, and should be reported as such.

\subsection{Worked reading}

On this project's validation years (targets 2015--2018, 153 countries),
the linear-trend baseline has point skill $-0.078$ against persistence, but
the country-clustered 95 percent interval (Equation~\ref{eq:boot-ci}) spans
roughly $[-0.34, +0.20]$ with a bootstrap $p$-value near $0.7$. The honest
conclusion is therefore not ``trend loses to persistence'' but ``trend and
persistence are statistically indistinguishable on this sample'': the point
gap is well within sampling noise once the country-level dependence is
accounted for. Reporting the interval, not only the point estimate, is what
keeps the evaluation honest at the level of statistical significance, which
is the same discipline this project already applies at the level of the
baseline comparison itself.


\bibliographystyle{plain}
\bibliography{references}

\end{document}
